
\chapter{Metodología}
\label{cap:metodologia}

Este capítulo describe el enfoque metodológico que vertebra el trabajo. El
problema abordado consiste en proponer una transformación de código, construir
un prototipo que la realice y comprobar empíricamente sus efectos; esa
naturaleza encaja de forma natural con el paradigma de \emph{Design Science
Research} (DSR), orientado a generar conocimiento mediante la construcción y
evaluación de artefactos.

\section{Design Science Research}
\label{sec:dsr-metodo}

DSR aborda la investigación a través de un artefacto que resuelve un problema
relevante y cuya utilidad se demuestra mediante evaluación rigurosa. Hevner et
al.~\cite{10.2307/25148625} sitúan en el núcleo de este paradigma la
articulación entre construir el artefacto y evaluarlo frente al problema que
pretende resolver. Peffers et al.~\cite{10.2753/MIS0742-1222240302} concretan
ese principio en una secuencia de seis actividades: identificar el problema y
motivarlo, fijar los objetivos de la solución, diseñar y desarrollar el
artefacto, demostrarlo, evaluarlo y comunicarlo. Esta secuencia proporciona una
estructura reconocible para presentar este tipo de investigación. La
\autoref{fig:dsr-ciclo} representa ese ciclo y la realimentación entre
construcción y evaluación que lo caracteriza.

% TODO: REGENERAR figuras/img/DSR.png desde el fuente actualizado
% figuras/src/fig-dsr-ciclo.mmd (el texto "impacto en CC" se cambió a
% "impacto en complejidad cognitiva"; el PNG actual aún muestra "CC").
\begin{figure}[ht]
\centering
\includegraphics[width=0.4\linewidth]{figuras/img/DSR.png}
\caption{Ciclo de \emph{Design Science Research} seguido en el trabajo, según
las actividades de Peffers et al.~\cite{10.2753/MIS0742-1222240302}. La
construcción del artefacto y su evaluación se articulan en un ciclo con
realimentación; las anotaciones indican su materialización en este TFM.}
\label{fig:dsr-ciclo}
\end{figure}

El artefacto de este trabajo es un prototipo determinista de detección y
refactorización de condicionales anidados. Su carácter determinista cumple
además una función metodológica de segundo orden: actúa como \emph{baseline}
verificable, es decir, como la referencia fiable contra la que se contrastan
después las refactorizaciones producidas por los grandes modelos de lenguaje
en RQ3. De este modo, la doble naturaleza de construir y comprobar, propia de
DSR, no solo organiza la propuesta, sino que sostiene el diseño experimental.



\section{Operacionalización en el TFM}
\label{sec:dsr-operacion}
\label{sec:operacionalizacion}

La evaluación del artefacto se articula en torno a las tres preguntas de
investigación: qué casos de condicionales anidados son combinables y con qué
resultado (RQ1), qué efecto tiene la transformación sobre la complejidad
cognitiva (RQ2) y en qué medida un modelo de lenguaje reproduce la
transformación de forma correcta y automática (RQ3). La
\autoref{tab:dsrm-tfm} relaciona las actividades de la metodología de Peffers
et al. con su materialización concreta en este trabajo y con los capítulos que
las desarrollan.

\begin{table}[ht]
\centering
\caption{Actividades de la metodología DSR de Peffers et
al.~\cite{10.2753/MIS0742-1222240302} y su materialización en el trabajo.}
\label{tab:dsrm-tfm}
\begin{tabular}{@{}p{4.2cm}p{5.2cm}p{2.6cm}@{}}
\toprule
\textbf{Actividad (Peffers et al.)} & \textbf{Materialización en el trabajo} & \textbf{Capítulo} \\
\midrule
Identificar el problema y motivarlo & Coste cognitivo del anidamiento de condicionales & \ref{cap:problema} \\
Fijar los objetivos de la solución  & Objetivos y preguntas de investigación RQ1--RQ3 & \ref{cap:introduccion} \\
Diseñar y desarrollar el artefacto  & Precondiciones P1--P4, prototipo basado en AST y oráculo & \ref{cap:propuesta} \\
Demostrar el artefacto              & Ejecución sobre el corpus piloto y real & \ref{sec:experimento} \\
Evaluar                             & Medición de $\Delta$ de complejidad y evaluación de los LLMs & \ref{sec:resultados} \\
Comunicar                           & La presente memoria y el paquete de replicación & --- \\
\bottomrule
\end{tabular}
\end{table}

La parte de \emph{construcción} comprende la formalización de las
precondiciones de aplicabilidad y el desarrollo del prototipo
(\autoref{cap:propuesta}). La parte de \emph{evaluación} se apoya en
protocolos deliberadamente congelados (con todos sus parámetros fijados de
antemano y registrados, de modo que no varíen entre ejecuciones) y en
artefactos persistentes y auditables, de modo que el procedimiento pueda
repetirse y conducir al mismo resultado; sus detalles se describen en el diseño experimental
(\autoref{sec:experimento}).

Con ese marco, cada pregunta se operacionaliza, sin alterar su enunciado
(\autoref{sec:rq-intro}), en los términos siguientes.

\paragraph{RQ1 — \textquestiondown{}En qué casos se pueden combinar sentencias condicionales anidadas y cuál sería el resultado?}
Se aborda de forma constructiva: se delimita el patrón objetivo y se formalizan
las precondiciones que hacen segura su combinación, que el detector aplica sobre
el corpus para separar los casos elegibles de los que no lo son. En concreto:
\begin{itemize}
  \item Se estudian condicionales anidados de tipo \lstinline|if| en Java.
  \item Foco inicial: un \lstinline|if| que contiene exactamente otro
  \lstinline|if|, sin ramas \lstinline|else| ni \lstinline|else if|, y sin
  sentencias adicionales en el bloque externo.
  \item El resultado es la transformación
  \lstinline|if (A) { if (B) { S } }| $\rightarrow$ \lstinline|if (A && B) { S }|.
  \item Se documentan las precondiciones de aplicabilidad
  (\autoref{sec:precondiciones}).
\end{itemize}

\paragraph{RQ2 — \textquestiondown{}Qué impacto tienen estas refactorizaciones en la complejidad cognitiva?}
Se aborda de forma empírica: se mide la complejidad cognitiva de cada método
antes y después de aplicar la transformación y se agrega el efecto sobre todo el
corpus. En concreto:
\begin{itemize}
  \item La complejidad cognitiva se mide siguiendo el modelo de
  SonarSource~\cite{10.1145/3194164.3194186}, mediante una implementación
  \emph{proxy} cuyas limitaciones se declaran en la \autoref{sec:proxy}.
  \item Se compara el valor antes y después de la refactorización en cada caso.
  \item Se reporta la reducción absoluta y relativa, por caso y agregada.
\end{itemize}

\paragraph{RQ3 — \textquestiondown{}Pueden los grandes modelos de lenguaje realizar esta refactorización de forma correcta y automática?}
Se aborda de forma comparativa: se solicita la misma transformación a varios
modelos de lenguaje bajo un protocolo controlado y sus salidas se contrastan con
el \emph{baseline} determinista mediante un oráculo automático. En concreto:
\begin{itemize}
  \item Se evalúa sobre el corpus de evaluación (\autoref{sec:corpus}).
  \item Se solicita a uno o más LLMs la misma refactorización bajo un
  \emph{prompt} estandarizado y un protocolo congelado.
  \item Se comparan los resultados en corrección semántica, completitud de
  detección y calidad del código, mediante un oráculo automático
  (\autoref{fig:oraculo}).
\end{itemize}

\paragraph{Criterio de corrección.}
Una transformación se considera correcta si preserva el comportamiento
observable del método. Esa corrección se sostiene en dos niveles. Por
construcción: bajo las precondiciones de aplicabilidad
(\autoref{sec:precondiciones}), el operador \lstinline|&&| respeta el orden y
el cortocircuito de la evaluación original, de modo que ambas versiones
producen el mismo efecto para toda entrada. Y de forma dinámica: el resultado
de cada caso elegible se compila con el compilador de Java para comprobar que
la transformación no genera código inválido. La ejecución de las pruebas de
comportamiento de los proyectos de origen excede el alcance del prototipo,
pues requeriría reconstruir y configurar cada proyecto, y se señala como
trabajo futuro (\autoref{cap:conclusiones}).

\section{Hipótesis de trabajo}
\label{sec:hipotesis}
El trabajo se guía por tres hipótesis que se contrastan con los resultados y
que, deliberadamente, no anticipan cifras.

\begin{description}
  \item[H1.] Existe un subconjunto no trivial de condicionales anidados en
  código Java real para los que la combinación es segura bajo las
  precondiciones de aplicabilidad.
  \item[H2.] La combinación reduce la complejidad cognitiva (proxy) de forma
  medible en los casos elegibles.
  \item[H3.] Los LLMs pueden reproducir la transformación con tasas de acierto
  elevadas en casos elegibles, pero pueden fallar en el rechazo de casos no
  elegibles.
\end{description}
